% Planned detached-authentication experiment; requires amsmath.
% Acquisition metadata and final results remain pending.
\section{Trace-Driven Experimental Methodology}
\label{sec:methodology}

The experiment evaluates the cost of adding real digital signatures to
recorded ADS-B traffic. Ordinary messages remain eligible for transmission
independently of signing. Their reception and subsequent authentication are
measured separately, together with processing demand and additional
channel interference.

\subsection{Dataset}
\label{subsec:opensky-dataset}

The planned evaluation uses two one-hour observation windows from
[REGION, DATES, AND ACQUISITION SYSTEM], selected to represent different
recorded traffic intensities. These are empirical scenarios, not statistical
repetitions. OpenSky observations provide the target ADS-B messages
\cite{Schafer2014OpenSky,Strohmeier2014OpenSky}. Preprocessing retains valid
DF17 observations, their original 14-byte messages, aircraft addresses, and
timestamps. Repeated message contents are preserved: identical bytes alone
do not establish duplicate observations of one physical transmission.

An aligned channel trace supplies all recorded 1090~MHz events, including
non-target traffic, with their timing and justified durations. Undecodable
activity is included only if recorded by the acquisition system. A
decoded-message feed cannot establish complete RF occupancy; acquisition
coverage and limitations are reported. Each target links to one channel
event to avoid double counting. Every configuration uses the same input
for its window, without synthetic ordinary traffic. Final dataset counts
and acquisition metadata remain to be completed.

\subsection{Signing and Transmission}
\label{subsec:replay-simulator}

We compare ECDSA-P256, ML-DSA-44, Falcon-512, and SLH-DSA-SHA2-128s using
groups of $k\in\{1,5,10,20\}$ consecutive messages per aircraft. Groups close
only after exactly $k$ messages are available, without a timeout. Each
ordinary message becomes eligible for transmission at its trace timestamp,
while a copy is retained for signing. Complete groups enter the signing
queue. Incomplete final groups remain unsigned, but their messages still
transmit. Thus, $k=1$ is per-message detached authentication, without
collection delay, rather than instantaneous authentication.

Each aircraft has its own key pair within each signed workload. Real
signatures cover the ordered original bytes and explicitly encoded
association context. Signed objects are reused across channel and rate
configurations. The context contains aircraft, algorithm, session, sequence,
source-time, message-count, key-fingerprint, and ordered message-tag fields;
the repository specifies their byte layout.

For an $s$-byte signature, the stream contains $52+8k$ context bytes, two
signature-length bytes, the signature, and two stream-length bytes. Each
hypothetical seven-byte ME field carries a two-byte group sequence, a
two-byte fragment index, and three stream bytes. Therefore,
\[
 F_{\mathrm{prototype}}=\left\lceil\frac{56+8k+s}{3}\right\rceil,
 \qquad F_{\min}=\left\lceil\frac{s}{7}\right\rceil.
\]
The second quantity is an analytical signature-only lower bound omitting
metadata, not a second executable protocol. Fragments occupy
$120~\mu$s frames. This research encoding has no allocated operational
ADS-B message type. Keys and sessions are provisioned separately;
credential traffic, retransmission, and error correction are excluded.

Each aircraft has one simulated transmitter. Ready ordinary messages have
priority, but an ongoing authentication frame finishes before another
starts. Fragment starts are paced at $r\in\{10,50,100\}$ frames/s per
aircraft. This changes when traffic is offered, not the bytes required for
a fixed workload; increased collisions may offset faster transmission.
Non-target events cause receiver interference but do not reserve the
modeled aircraft transmitter.

\subsection{Computational Model}
\label{subsec:computational-model}

Host execution generates valid cryptographic objects; published embedded
measurements determine simulated processing times. ECDSA uses the
P256-Cortex-M4 report for an nRF52840 at 64~MHz \cite{P256CortexM4}.
Post-quantum timings use pqm4 cycle counts and its 24~MHz benchmark clock
\cite{PQM4Benchmarks,PQM4Platform}. Provisional FN-DSA and SPHINCS+
measurements approximate Falcon and SLH-DSA, respectively. ECDSA timings
exclude external hashing and nonce generation. Fixed representative times
omit execution-time and message-length variability. These heterogeneous
references define service-demand scenarios, not a controlled algorithm
ranking or measured aircraft performance.

There is one signing worker per aircraft and one receiver verification
worker. If verification queues grow persistently, affected configurations
are additionally evaluated with two and four receiver workers. Service
demand and backlog assess processing feasibility; memory, energy, and
contention with other applications require separate measurements.

\subsection{Channel Comparison}
\label{subsec:experimental-variables}

Each configuration is evaluated with a collision-free control and a
destructive-overlap model that discards overlapping frames within one
collision domain. The latter is an uncalibrated interference sensitivity
model, omitting received power, capture effects, and receiver geometry.
No artificial random erasure or timing jitter is added. Fixed timings and
cached signatures make each replay deterministic, so repeated random seeds
are unnecessary. The main matrix contains 96 configurations per window:
four algorithms, four group sizes, three rates, and two channel models.

The unsigned baseline retains recorded event times. The augmented replay
replaces each target event with its scheduled transmission and adds
fragments; other recorded events retain their times. A timing-matched
control removes fragments while retaining the augmented ordinary schedule.
These comparisons distinguish the overall authentication effect from
direct fragment interference. Inter-aircraft overlap causes losses, rather
than a channel-wide transmission waiting queue.

\subsection{Metrics and Observation Period}
\label{subsec:evaluation-metrics}

Let $t_0$, $t_{\mathrm{tx}}$, $t_{\mathrm{rx}}$, and $t_{\mathrm{auth}}$
denote a message's trace, transmission, reception, and successful
authentication times. We report
\[
 D_{\mathrm{tx}}=t_{\mathrm{tx}}-t_0,
 \qquad D_{\mathrm{auth}}=t_{\mathrm{auth}}-t_{\mathrm{rx}}.
\]
Collection and signing waits contribute to authentication latency; ordinary
delay reflects transmitter scheduling. Trace timestamps do not reveal
sensor measurement times, so these are added modeled delays rather than
total sensor-to-transmission latency.

The receiver retains ordinary messages, reconstructs delivered authentication
objects, and performs real verification using received bytes. Association
uses signed tags and source-time intervals with a conservative timing bound
derived from the input arrivals and transmitter rules; a shared clock is
assumed. All originals in a group are required. Results distinguish
ordinary reception, object reconstruction, verification attempts and
failures, successful authentication, and unfinished work. We report
authentication coverage over all targets and over received targets,
completed-case latency distributions, reception gaps, additional frames
and airtime, processor demand, and queue backlog. Offered airtime can
exceed elapsed time because transmissions overlap; it is not RF occupancy.

A fixed 600~s follow-up is declared before final runs. Queues start empty
at each window's beginning. An isolated largest
SLH-DSA group at $r=10$ requires approximately 589~s for signing, fragment
delivery, and verification under the selected profiles, motivating this
observation budget. Queue clearance is not guaranteed. No new target
groups begin during follow-up, while recorded surrounding traffic
continues through the cutoff, including frame completion. Unfinished work
remains pending rather than a definitive failure. Delay-threshold coverage
includes only messages with sufficient observed follow-up, with exclusions
reported.

\subsection{Reproducibility}
\label{subsec:reproducibility}

Inputs, source versions, benchmark references, configurations, and checksums
are retained. Cached signatures and public keys preserve the exact objects
across replays; regenerating them may change their bytes or lengths.
Separate commands and output directories reproduce both windows. The
existing sample validates implementation; acquisition, final execution,
and analysis of the two windows remain pending.
